% GENERATED FILE. Edit canonical lessons, then run npm run generate:books.
% canonical-source-hash: fnv1a64:1eb34db84ca73673
% canonical-lessons: ZH-C111-qianbian, ZH-W111-bian, ZH-C111-houbian, ZH-C111-fujin, ZH-C111-lukou, ZH-C111-guai

\chapter{In Front and Behind}
\label{ch:zh-in-front-and-behind}
\hlchaptermodality{111}

\begin{chapteropening}
\textbf{By the end of this chapter:} \emph{I can say in front, behind, nearby and at the crossing, and say where to turn.}

Complete the last lesson of chapter 111: I can say in front, behind, nearby and at the crossing, and say where to turn.
\end{chapteropening}

\section[\texorpdfstring{\zh{前边} (qiánbian)}{前边 (qiánbian)}]{\zh{前边} (qiánbian) — in front; ahead}

\label{lesson:ZH-C111-qianbian}



\begin{quote}
Before the new one: say the Mandarin for to paint pictures, then the Mandarin for to dance.
\end{quote}

\subsection*{You'll want to know: \zh{前边}}
\textbf{\zh{前边}} — \emph{qiánbian} — \textquotedblleft{}in front; ahead\textquotedblright{}.

\begin{quote}
\textbf{\zh{前边是火车站}} \emph{qiánbian shì huǒchēzhàn} — The railway station is up ahead.
\end{quote}

The new character, in pieces: \zh{边} — lì, the foot of \zh{男}, inside the walking piece that runs down the left and along the foot of \zh{这} and \zh{进}. The next lesson puts it in your hand.

\begin{cousinweb}
\textbf{\zh{前边}} (\emph{qiánbian}) is built from \zh{前} 'front' and \zh{边} 'side', both inherited from Middle Chinese, \emph{dzen} and \emph{pen} (Guangyun, in Baxter's transcription); the dictionary glosses the first 'first, before' and the second 'an edge'.
\end{cousinweb}

\begin{grammarlens}[title={What you've built}]
One more word for asking where things are.
\end{grammarlens}

\subsection*{Your turn}
\begin{itemize}
\raggedright
  \item \emph{Say it:} \emph{qiánbian}
  \item \emph{Say it:} \emph{qiánbian}, once more
  \item \emph{Say it:} \emph{qiánbian shì huǒchēzhàn} — The railway station is up ahead.
  \item \emph{Recall:} say the Mandarin for to paint pictures, then the Mandarin for to dance, then say \emph{qiánbian} again
\end{itemize}

\begin{tcolorbox}[breakable,colback=teal!4,colframe=teal!35!black,title={Before you move on}]
What does \zh{前边} mean? (\textquotedblleft{}In front; ahead\textquotedblright{}.) Say it once more.
\end{tcolorbox}
\section[\texorpdfstring{\zh{边} (biān)}{边 (biān)}]{\zh{边} — the foot of \zh{男}, inside the walking piece}

\label{lesson:ZH-W111-bian}



\begin{quote}
The word first, then the character inside it.

\begin{itemize}
\raggedright
  \item \emph{Recall:} say \emph{qiánbian} — \textbf{R1}, one lesson back, before it decays
\end{itemize}

\begin{quote}
\zh{男} \zh{这} \zh{进}
\end{quote}

\emph{Read it:} each one aloud, with nothing said first — \emph{nán}, \emph{zhè}, \emph{jìn}
\end{quote}

\subsection*{Script — the foot of \zh{男}, inside the walking piece}
\hlblockfigure{\detokenize{figures/ZH-W111-bian-filmstrip.pdf}}{How \zh{边} is written, stroke by stroke}

\begin{quote}
\zh{边}
\end{quote}

\textbf{The inside first, two strokes}, the foot of \zh{男}: a stroke across that turns down and hooks, then a left-falling stroke through it.

\textbf{The walking piece, three strokes}, written last, as in \zh{这} and \zh{进}: a dot, a small zigzag down the left, and a long flat sweep along the foot.

2 + 3 = five.

\subsection*{Your turn}
\begin{itemize}
\raggedright
  \item \emph{Write it:} \zh{边}, all five, with the model in front of you
  \item \emph{Cover:} write it again from memory, and count the pieces, not the strokes
\end{itemize}

\begin{tcolorbox}[breakable,colback=teal!4,colframe=teal!35!black,title={Before you move on}]
Which comes first in \zh{边}, the inside or the walking piece? (\textbf{The inside}.) Where have you written that inside piece? (\textbf{At the foot of \zh{男}}.) Which characters you know end with the walking piece? (\textbf{\zh{这}} and \textbf{\zh{进}}.)

Source: \href{https://raw.githubusercontent.com/chanind/hanzi-writer-data/68d10a4b21150cae5e1ebbd223eed289cf32d90c/data/\%E8\%BE\%B9.json}{Hanzi Writer Data \zh{边}, pinned stroke paths and medians}.
\end{tcolorbox}
\section[hòubian]{hòubian — behind; at the back}

\label{lesson:ZH-C111-houbian}



\begin{quote}
Before the new one: say the Mandarin for to dance, then the Mandarin for in front.
\end{quote}

\subsection*{You'll want to know: hòubian}
\emph{hòubian} — \textquotedblleft{}behind; at the back\textquotedblright{}.

\begin{quote}
\emph{xuéxiào hòubian yǒu yí ge gōngyuán} — There's a park behind the school.
\end{quote}

Its characters come later; for now it is a word for the ear.

\begin{cousinweb}
\emph{hòubian} is built from hòu 'behind' and \zh{边} 'side', both inherited from Middle Chinese, \emph{huwX} and \emph{pen} (Guangyun, in Baxter's transcription).
\end{cousinweb}

\begin{grammarlens}[title={What you've built}]
One more word for asking where things are.
\end{grammarlens}

\subsection*{Your turn}
\begin{itemize}
\raggedright
  \item \emph{Say it:} \emph{hòubian}
  \item \emph{Say it:} \emph{hòubian}, once more
  \item \emph{Say it:} \emph{xuéxiào hòubian yǒu yí ge gōngyuán} — There's a park behind the school.
  \item \emph{Recall:} say the Mandarin for to dance, then the Mandarin for in front, then say \emph{hòubian} again
\end{itemize}

\begin{tcolorbox}[breakable,colback=teal!4,colframe=teal!35!black,title={Before you move on}]
What does hòubian mean? (\textquotedblleft{}Behind; at the back\textquotedblright{}.) Say it once more.
\end{tcolorbox}
\section[fùjìn]{fùjìn — nearby; in the area}

\label{lesson:ZH-C111-fujin}



\begin{quote}
Before the new one: say the Mandarin for in front, then the Mandarin for at the back.
\end{quote}

\subsection*{You'll want to know: fùjìn}
\emph{fùjìn} — \textquotedblleft{}nearby; in the area\textquotedblright{}.

\begin{quote}
\emph{fùjìn yǒu yínháng ma} — Is there a bank near here?
\end{quote}

Its characters come later; for now it is a word for the ear.

\begin{cousinweb}
\emph{fùjìn} is built from fù 'to attach' and jìn 'near', both inherited from Middle Chinese, \emph{bjuH} and \emph{gj+nX} (Guangyun, in Baxter's transcription).
\end{cousinweb}

\begin{grammarlens}[title={What you've built}]
One more word for asking where things are.
\end{grammarlens}

\subsection*{Your turn}
\begin{itemize}
\raggedright
  \item \emph{Say it:} \emph{fùjìn}
  \item \emph{Say it:} \emph{fùjìn}, once more
  \item \emph{Say it:} \emph{fùjìn yǒu yínháng ma} — Is there a bank near here?
  \item \emph{Recall:} say the Mandarin for in front, then the Mandarin for at the back, then say \emph{fùjìn} again
\end{itemize}

\begin{tcolorbox}[breakable,colback=teal!4,colframe=teal!35!black,title={Before you move on}]
What does fùjìn mean? (\textquotedblleft{}Nearby; in the area\textquotedblright{}.) Say it once more.
\end{tcolorbox}
\section[lùkǒu]{lùkǒu — crossroads; junction}

\label{lesson:ZH-C111-lukou}



\begin{quote}
Before the new one: say the Mandarin for at the back, then the Mandarin for nearby.
\end{quote}

\subsection*{You'll want to know: lùkǒu}
\emph{lùkǒu} — \textquotedblleft{}crossroads; junction\textquotedblright{}.

\begin{quote}
\emph{yínháng zài lùkǒu} — The bank is at the crossroads.
\end{quote}

Its characters come later; for now it is a word for the ear.

\begin{cousinweb}
\emph{lùkǒu} is built from lù 'road' and \zh{口} 'mouth', both inherited from Middle Chinese, \emph{luH} and \emph{khuwX} (Guangyun, in Baxter's transcription): the mouth of a road.
\end{cousinweb}

\begin{grammarlens}[title={What you've built}]
One more word for asking where things are.
\end{grammarlens}

\subsection*{Your turn}
\begin{itemize}
\raggedright
  \item \emph{Say it:} \emph{lùkǒu}
  \item \emph{Say it:} \emph{lùkǒu}, once more
  \item \emph{Say it:} \emph{yínháng zài lùkǒu} — The bank is at the crossroads.
  \item \emph{Recall:} say the Mandarin for at the back, then the Mandarin for nearby, then say \emph{lùkǒu} again
\end{itemize}

\begin{tcolorbox}[breakable,colback=teal!4,colframe=teal!35!black,title={Before you move on}]
What does lùkǒu mean? (\textquotedblleft{}Crossroads; junction\textquotedblright{}.) Say it once more.
\end{tcolorbox}
\section[guǎi]{guǎi — to turn a corner}

\label{lesson:ZH-C111-guai}



\begin{quote}
Before the new one: say the Mandarin for nearby, then the Mandarin for crossroads.
\end{quote}

\subsection*{You'll want to know: guǎi}
\emph{guǎi} — \textquotedblleft{}to turn a corner\textquotedblright{}.

\begin{quote}
\emph{zhèr néng guǎi ma} — Can we turn here?
\end{quote}

Its characters come later; for now it is a word for the ear.

\begin{cousinweb}
\emph{guǎi} is inherited from Middle Chinese \emph{kweaX} (Guangyun, in Baxter's transcription), glossed as a stick, a thing for the hands and feet, probably a crutch; how it came to mean turning a corner is not settled.
\end{cousinweb}

\begin{grammarlens}[title={What you've built}]
One more word for asking where things are.
\end{grammarlens}

\subsection*{Your turn}
\begin{itemize}
\raggedright
  \item \emph{Say it:} \emph{guǎi}
  \item \emph{Say it:} \emph{guǎi}, once more
  \item \emph{Say it:} \emph{zhèr néng guǎi ma} — Can we turn here?
  \item \emph{Recall:} say the Mandarin for nearby, then the Mandarin for crossroads, then say \emph{guǎi} again
\end{itemize}

\begin{tcolorbox}[breakable,colback=teal!4,colframe=teal!35!black,title={Before you move on}]
What does guǎi mean? (\textquotedblleft{}To turn a corner\textquotedblright{}.) Say it once more.
\end{tcolorbox}
